% paper/sections/appendix_verification_details.tex
%
% 付録 G: 検証の補足解析
%   G.1 PPE 離散化の選択（§12.5 から移動）
%   G.2 HFE の smoothed-Heaviside 圧力場への適合性（§12.4 から移動）
%   G.3 V7 IMEX-BDF2 二相時間精度 RCA（§13.4 補足）
% ═══════════════════════════════════════════════════════════════════════════════
\section{検証の補足解析}
\label{sec:appendix_verification}

% ─────────────────────────────────────────────────────────────────────────────
\subsection{PPE 離散化の選択：FD 直接解法と CCD 勾配の併用}
\label{sec:ppe_discretization_choice}
% ─────────────────────────────────────────────────────────────────────────────

§~\ref{sec:interface_crossing} で smoothed-Heaviside 一括解法の密度比依存性を分析したが，
CCD PPE の欠陥補正（DC）にはさらに\textbf{本質的制約}がある：
CCD 高次演算子 $L_H$ と低次前処理 $L_L$（2 次精度 FD）の
\textbf{スペクトル不整合}により，界面型密度場に限らず定常反復が発散する．

\paragraph{反復行列のスペクトル解析}

DC 反復 $p^{k+1} = p^k + L_L^{-1}(q_h - L_H p^k)$ の収束条件は
反復行列 $T = I - L_L^{-1} L_H$ のスペクトル半径 $\rho(T) < 1$ である．
$N = 16$, $\rho_l/\rho_g = 2$ の静止液滴問題で $T$ の全固有値を計算した結果：
\begin{itemize}[noitemsep]
  \item $L_L^{-1} L_H$ の固有値実部範囲：$[-0.03,\; 2.45]$（\textbf{不定}）
  \item $\rho(T) = 1.45 > 1$
  \item $|\lambda| > 1$ の固有値：289 自由度中 40 個（14\%）
\end{itemize}
$L_L^{-1} L_H$ が\textbf{不定}（負の固有値を持つ）であるため，
緩和係数 $\omega$ の最適化を行っても $\rho(I - \omega L_L^{-1} L_H) \approx 1$ となり，
実用的な収束は得られない．

\paragraph{不定性の原因}

CCD コンパクトスキームの修正波数 $k^*_{\mathrm{CCD}}$ は
高波数領域で真の波数 $k$ を\textbf{超える}（super-resolution 特性）．
一方 2 次精度 FD の修正波数 $k_{\mathrm{FD}}$ は全波数で $k$ 以下である．
この比 $k^*_{\mathrm{CCD}} / k_{\mathrm{FD}}$ が高波数で 1 を超えるため，
$L_L^{-1} L_H$ の固有値が 1 を超え，
Richardson 型反復（DC, Jacobi, GS, ADI を含む）が発散する．
この不整合は密度比とは独立であり，\textbf{単相でも発生する}．

\paragraph{代替反復法の系統的検証}

DC の収束を回復する 7 つの方法を検証した：

\begin{table}[htbp]
\centering
\caption{CCD PPE 反復ソルバの代替法と結果}
\label{tab:dc_alternatives}
\begin{tabular}{>{\raggedright\arraybackslash}p{40mm}cl}
\toprule
方法 & 残差到達値 & 判定 \\
\midrule
DC（フィルタなし）        & $\sim 10^{-2} \to$ 発散  & 不可 \\
DC + 低域フィルタ          & $\sim 10^{-3}$（停滞）   & 精度不足 \\
DC + フィルタ漸減          & 漸減解除時に再発散        & 不可 \\
CCD 1D ライン緩和          & 発散（ADI 構造と同根）   & 不可 \\
HFE 適用（残差/補正/解）   & 悪化（滑らかな場に逆効果）& 不可 \\
DC + 早期停止               & $\sim 10^{0}$（安定だが粗い）& $22\times$ 劣化 \\
Anderson 加速（$m{=}30$）  & $\sim 10^{-2}$（収束遅い）& 反復数過大 \\
\bottomrule
\end{tabular}
\end{table}

\paragraph{FD 直接解法 + CCD 勾配：投影不整合の評価}

CCD PPE ($L_H p = q_h$) を安価な反復法で解くことは現状では不可能である．
代替として，\textbf{2 次精度 FD の PPE を直接法で厳密に解き}，
速度補正の勾配 $\bnabla p$ には\textbf{CCD を適用}する：
\begin{alignat}{2}
  &\text{PPE：}  &\quad L_{\mathrm{FD}}\, p &= q_h
    \qquad\text{（FD 直接法，残差 $= 0$）} \label{eq:fd_ppe}\\
  &\text{速度補正：} &\quad \bm{u}^{n+1} &= \bm{u}^* - \frac{\Delta t}{\tilde\rho}\,
    D^{(1)}_{\mathrm{CCD}} p
    \qquad\text{（CCD $\Ord{h^6}$）} \label{eq:ccd_correction}
\end{alignat}
投影不整合は $\Ord{h^2}$ だが，各ステップで 1 回だけ生じるため有界に留まる．
表面張力 $\bm{f}_\sigma$ と圧力勾配が\textbf{同一の CCD 演算子}で評価されるため，
Balanced-Force 条件は完全に維持される．

\begin{table}[htbp]
\centering
\caption{PPE 離散化方式の比較}
\label{tab:ppe_method_comparison}
\PaperCaptionNote{静止液滴，$N = 64$，200 ステップで比較する．}
\begin{tabular}{lccl}
\toprule
方式 & $\|\bm{u}\|_\infty$ & $\Delta p$ 誤差 & 備考 \\
\midrule
FD PPE + CCD $\bnabla p$ & $2.3 \times 10^{-4}$ & 1.2\% & \textbf{採用} \\
FD PPE + FD $\bnabla p$  & $1.1 \times 10^{-2}$ & 1.0\% & BF 崩壊 \\
DC 早期停止 + CCD $\bnabla p$ & $5.1 \times 10^{-3}$ & 1.1\% & 残差蓄積 \\
\bottomrule
\end{tabular}
\end{table}

% ─────────────────────────────────────────────────────────────────────────────
\subsection{HFE の smoothed-Heaviside 圧力場への適合性}
\label{sec:hfe_smoothed_scope}
% ─────────────────────────────────────────────────────────────────────────────

§~\ref{sec:field_extension} で導入した HFE（Hermite 場延長法）は
IPC における $\bnabla p^n$ の界面ジャンプ処理のために設計されたが，
smoothed-Heaviside 解法では圧力場が $\varepsilon$-正則化により
\textbf{既に滑らか}であるため，HFE は逆効果となる．

静止液滴テスト（$R=0.25$，$\mathit{We} = 10$，非増分投影法，200 ステップ）で
HFE ON/OFF を比較した結果，HFE を適用すると寄生流れが 4 桁以上増大し
Laplace 圧精度が完全に破綻した．
これは HFE がターゲット相の圧力値をソース相の値で上書きすることで
界面近傍に人為的な勾配不連続が導入されることに起因する．

HFE の適用範囲は以下のように明確に規定される：
\begin{enumerate}[noitemsep]
  \item \textbf{smoothed-Heaviside + 非増分投影}：
        HFE は不要（有害）．圧力場は $\varepsilon$-正則化で滑らか．
  \item \textbf{増分投影法（IPC）}：
        $p^n$ は Young--Laplace ジャンプ $j_{gl}=p_g-p_l=-\sigma\kappa_{lg}$ を持つため，
        HFE による延長が\textbf{必須}．
  \item \textbf{分相 PPE 解法}：
        界面上のジャンプ条件を HFE で課す必要がある
        （§~\ref{sec:split_ppe_recovery}）．
\end{enumerate}

% ─────────────────────────────────────────────────────────────────────────────
\subsection{V7 IMEX-BDF2 二相時間精度の原因分析}
\label{app:v7_imex_bdf2_rca}
% ─────────────────────────────────────────────────────────────────────────────

本付録は §~\ref{sec:imex_bdf2_twophase_time} の V7 に対する追加 RCA
（root-cause analysis）を記録する．結論は，観測 slope $1.48$--$1.59$ が
IMEX-BDF2 単体の係数誤りではなく，
\textbf{capillary pressure-jump / affine FCCD projection が界面帯で作る低正則性・分裂誤差}
に支配される，というものである．
したがって V7 は BDF2 単体次数の試験ではなく，§~\ref{sec:validation} と同じ
capillary coupled stack の実効時間精度診断として読む．
BDF2 単体の設計次数は U8（§~\ref{sec:U8_time_integration}）で独立に検証済みである．

\paragraph{診断設計}

基準 run は $N=24$, $T=0.02$, $\rho_l/\rho_g=10$, $\sigma=1$ とし，
$n_{\mathrm{step}}=\{8,16,32,64,128\}$ の系列で
$n_{\mathrm{ref}}=128$ を参照解に用いた．
元の V7 と同じ各 step 再初期化に加え，再初期化回数固定，再初期化なし，
界面静止，表面張力除去，一様係数化を分離して，低次数の必要条件を反証した．

\begin{table}[htbp]
\centering\scriptsize
\caption{V7 RCA：対照実験による支配因子の切り分け}
\label{tab:v7_rca_controls}
\begin{tabularx}{\linewidth}{>{\raggedright\arraybackslash}p{34mm}r r X}
\toprule
対照 run & finest slope & $n=64$ velocity $\|e\|_\infty$ & 解釈 \\
\midrule
\texttt{each\_step} & $1.59$ & $3.666\times10^{-3}$ &
$n_{\mathrm{ref}}=64$ はやや粗いが，参照解を深くしても 2 次へは戻らない．\\
\texttt{fixed\_reinit\_count} & $1.50$ & $9.358\times10^{-4}$ &
再初期化 cadence は誤差定数を変えるが，指数の主因ではない．\\
\texttt{no\_reinit} & $1.57$ & $1.670\times10^{-3}$ &
再初期化そのものも主因ではない．\\
\texttt{static\_interface} & $1.58$ & $1.870\times10^{-3}$ &
界面輸送を止めても低次数が残るため，TVD-RK3/IMEX-BDF2 の Lie 分割は必要条件ではない．\\
\texttt{static\_no\_capillary} & $1.54$ & $6.831\times10^{-8}$ &
$\sigma=0$ で誤差振幅が 4--5 桁崩壊し，capillary jump が支配因子であることを示す．\\
\texttt{uniform\_static\_capillary} & $1.59$ & $8.198\times10^{-3}$ &
$\rho,\mu$ を一様化しても $\sigma=1$ なら大誤差が残り，密度・粘性 jump は振幅源ではない．\\
一様係数，静止，$\sigma=0$ & $1.58$ & $7.302\times10^{-8}$ &
一様係数かつ $\sigma=0$ では near-zero error となる．\\
\bottomrule
\end{tabularx}
\end{table}

\paragraph{仮説監査}

表~\ref{tab:v7_rca_hypotheses} は，物理・数学モデルから自然に生じる候補原因を
支配因子，副次因子，反証済みに分類したものである．ここで「反証」は
該当メカニズムが完全にゼロであることではなく，V7 の観測次数低下を説明する
一次支配因子ではないことを意味する．

\begin{table}[htbp]
\centering\scriptsize
\caption{V7 RCA：仮説監査}
\label{tab:v7_rca_hypotheses}
\begin{tabularx}{\linewidth}{>{\raggedright\arraybackslash}p{9mm}X>{\raggedright\arraybackslash}p{23mm}X}
\toprule
ID & 仮説 & 判定 & 根拠 \\
\midrule
H01 & $n_{\mathrm{ref}}=64$ が粗すぎる & 副次因子 &
ref128 で slope は $1.48\to1.59$ に改善するが，2 次には戻らない．\\
H02 & step-based reinit cadence が主因 & 反証 &
全系列で reinit 回数を固定しても slope $1.50$．\\
H03 & 再初期化そのものが主因 & 反証 &
no-reinit でも slope $1.57$．\\
H04 & 界面輸送と IMEX の Lie split が必要条件 & 反証 &
static-interface でも slope $1.58$．\\
H05 & capillary pressure-jump projection が支配 & 支持 &
$\sigma=0$ で velocity error が $10^{-3}$--$10^{-2}$ から $10^{-7}$ へ崩壊する．\\
H06 & 密度比が振幅を支配 & 反証 &
一様係数でも $\sigma=1$ なら $8.198\times10^{-3}$ の誤差が残る．\\
H07 & 粘性比が振幅を支配 & 反証 &
H06 と同じ一様係数対照で否定される．\\
H08 & $\|e\|_\infty$ norm だけが判定を悪化 & 部分支持 &
$L^2$ でも同様の slope 傾向を示すが，最大誤差は界面局在である．\\
H09 & bulk velocity dynamics が律速 & 反証 &
bulk error は interface-band error より数桁小さい．\\
H10 & IMEX-BDF2 係数実装が誤り & 反証 &
U8 は BDF2 単体を通し，$\sigma=0$ 対照は near-zero error となる．\\
H11 & 旧 nodal pressure-history bug が残存 & 反証 &
現行離散化は projection-native face acceleration を使い，残差は capillary band に局在する．\\
H12 & BDF1 startup が主因 & 非支持 &
全対照で指数は近く，capillary 有無だけが振幅を桁で変える．\\
H13 & 空間格子 $N=24$ だけが問題 & 副次因子 &
同じ格子でも $\sigma=0$ では near-zero error となる．\\
H14 & PPE tolerance / DC iteration が主因 & 非支持 &
solver 設定を変えずに $\sigma=0$ とするだけで誤差が崩壊する．\\
H15 & 曲率・再初期化更新頻度が主因 & 反証 &
界面と曲率を固定しても capillary-limited behavior が残る．\\
H16 & 圧力不連続の低正則性が max-norm 時間収束を制限 & 支持 &
最大誤差は $\psi\simeq0.5$ に局在し，capillary jump 単独で大誤差を再現する．\\
H17 & 壁境界条件が支配 & 反証 &
最大誤差位置は壁ではなく円形界面上である．\\
H18 & volume/mass drift が速度 norm を汚染 & 反証 &
体積 drift は round-off floor に留まる．\\
H19 & GPU/CPU 経路差が原因 & 非支持 &
remote production path は本文値を再現し，局在構造も物理対照と整合する．\\
H20 & V7 を純粋 BDF2 次数検証として読むべき & 反証 &
V7 は capillary pressure-jump coupled-stack 診断であり，純粋 BDF2 は U8 が担う．\\
\bottomrule
\end{tabularx}
\end{table}

\paragraph{根本原因}

Young--Laplace pressure jump は界面上の分布的な不連続条件であり，
affine FCCD PPE では cut face 上の圧力 jump と速度補正が同じ投影過程で閉じる．
このとき bulk 解が滑らかでも，$\psi\simeq0.5$ の smeared interface band では
時間発展の正則性が bulk と同等にはならない．
実際，最大 velocity error は界面帯に局在し，bulk error と体積 drift は支配的でない．
したがって V7 の $\triangle$ 判定は，閾値の緩和や reinit tuning で合格化する対象ではなく，
capillary jump/projection の時間離散構造を上げない限り残る
Type-D の theoretical hard limit として扱う．
